\begin{figure}[H]
\centering
\begin{tikzpicture}
\begin{axis}[
    xbar, width=0.8\linewidth, height=11cm,
    title={\textbf{Stage 1 --- Per-Class Support (test set)}},
    xlabel={Number of positive samples},
    ytick=data, symbolic y coords={
        foreign-body
        normal-z-line
        normal-pylorus
        normal-cecum
        retroflex-rectum
        retroflex-stomach
        ileocecal-valve
        dyed-lifted-polyp
        dyed-resection-margins
        hemorrhoid
        diverticulum
        erosion
        gastritis
        esophagitis
        ulcerative-colitis
        crohns-disease
        barretts-esophagus
        gastric-ulcer
        duodenal-ulcer
        instrument
        polyp-pedunculated
        polyp-hyperplastic
        polyp-sessile
    },
    y dir=reverse,
    nodes near coords, nodes near coords style={font=\scriptsize},
    tick label style={font=\scriptsize}, label style={font=\small},
    title style={font=\small\bfseries}, grid=major, grid style={gray!20},
    bar width=5pt,
]
\addplot[fill=orange!70, draw=orange!80] coordinates {(3387,foreign-body) (3524,normal-z-line) (3524,normal-pylorus) (3524,normal-cecum) (3524,retroflex-rectum) (3524,retroflex-stomach) (3524,ileocecal-valve) (3546,dyed-lifted-polyp) (3546,dyed-resection-margins) (3586,hemorrhoid) (3586,diverticulum) (3635,erosion) (3635,gastritis) (3635,esophagitis) (3635,ulcerative-colitis) (3635,crohns-disease) (3635,barretts-esophagus) (3635,gastric-ulcer) (3635,duodenal-ulcer) (3681,instrument) (3709,polyp-pedunculated) (3709,polyp-hyperplastic) (3709,polyp-sessile)};
\end{axis}
\end{tikzpicture}
\caption[Stage 1 per-class support]{Per-class support (number of positive test
images per disease). Classes with very low support should be read alongside
their F1 in Figure~\ref{fig:stage1_per_class_f1}.}
\label{fig:stage1_support}
\end{figure}
